% ---------------------------------------------------------------------------
% Supplementary material -- Pseudo-code of the SA neighborhood.
% Paper: Parallel machine scheduling with simultaneous interruptions:
%        a case study in shoe sole production.
% Standalone document: compile with pdflatex.
% ---------------------------------------------------------------------------
\documentclass[11pt]{article}

\usepackage[utf8]{inputenc}
\usepackage{amsmath,amssymb}
\usepackage[linesnumbered,ruled,vlined]{algorithm2e}
\usepackage{geometry}
\usepackage[colorlinks=true,linkcolor=blue,urlcolor=blue]{hyperref}

\geometry{a4paper,margin=2.5cm}

\SetKwProg{Fn}{Function}{:}{end}
\SetKwFunction{Relocate}{Relocate}
\SetKwFunction{Redistribute}{Redistribute}
\SetKwFunction{TwoStep}{TwoStepNeighbor}
\SetKwFunction{Split}{SplitUniformly}

\title{Pseudo-code of the neighborhood generation\\ in the Simulated Annealing algorithm}
\author{Supplementary material to the paper\\
\emph{Parallel machine scheduling with simultaneous interruptions:}\\
\emph{a case study in shoe sole production}}
\date{}

\begin{document}
\maketitle

\section*{Solution representation}

A feasible solution is represented by a $p$-tuple $S=(S_1,\dots,S_p)$, one entry
per shelf, where each shelf
\[
  S_\ell = \bigl(\sigma_{\ell,1},\dots,\sigma_{\ell,|S_\ell|}\bigr)
\]
is the ordered sequence of the \emph{sub-jobs} processed on it, in the order in
which they are processed. A sub-job is a triple $\sigma=(j,i,q)$ formed by a job
$j\in J$ (a shoe size), one of its mold copies $i \in \{1,\dots,c_j\}$, and the
quantity $q\in\mathbb{Z}_{\geq 0}$ of soles produced with that copy. We write
\[
  \Sigma(S) = \bigl\{(\ell,r) : 1\leq \ell \leq p,\ 1 \leq r \leq |S_\ell| \bigr\}
\]
for the set of positions occupied by the sub-jobs of $S$, and $S_\ell[r]$ for the
sub-job in position $r$ of shelf $\ell$.

Every solution generated by the algorithm satisfies the two invariants
\begin{itemize}
  \item[(I1)] each pair $(j,i)$, with $j\in J$ and $i\in\{1,\dots,c_j\}$, occurs
    exactly once in $S$;
  \item[(I2)] $\sum_{i=1}^{c_j} q_{j,i} = n_j$ for every job $j \in J$, where
    $q_{j,i}$ is the quantity of the sub-job $(j,i)$ and $n_j$ is the number of
    soles of size $j$ required in the order.
\end{itemize}
The initial solution is built by splitting the quantity $n_j$ of each job over
its $c_j$ mold copies with the procedure \Split{} of Algorithm~\ref{alg:neigh},
and by appending each resulting sub-job to a shelf drawn uniformly at random,
jobs being taken in the order in which they appear in the instance and mold
copies in increasing order of $i$. Both operators described below preserve
(I1)--(I2), so every candidate solution generated during the search is feasible
and no repair step is required.

\section*{Neighborhood}

A neighbor of the current solution is generated by the procedure \TwoStep{} of
Algorithm~\ref{alg:neigh}, which combines two elementary moves:
\Relocate{}, a positional reallocation that moves a single sub-job to a randomly
chosen shelf and position without changing any quantity, and \Redistribute{},
a quantity redistribution that resamples how the quantity of a multi-mold job is
divided among its mold copies without changing any position. Three points of the
procedure deserve to be made explicit.

\begin{itemize}
  \item The sub-job drawn at the start of \TwoStep{} is used \emph{only} to
        decide which of the two moves is applied. Neither branch necessarily
        acts on it: \Relocate{} and \Redistribute{} draw their own sub-job,
        respectively their own job, independently of that first draw.
  \item The name ``two-step'' refers to the single-mold branch, in which two
        independent relocations are applied in sequence, the second one acting
        on the outcome of the first. The multi-mold branch applies a single
        redistribution and no relocation.
  \item In \Redistribute{}, the quantity $n_j$ of the selected job is
        partitioned uniformly at random over the vectors of $c_j$ non-negative
        integers summing to $n_j$, by the classical stars-and-bars scheme
        (procedure \Split{}). Zero parts are admissible, which allows the search
        to concentrate a job on a subset of its mold copies.
\end{itemize}

\begin{algorithm}[H]
  \caption{Neighbor generation in the SA algorithm}
  \label{alg:neigh}
  \KwIn{Solution $S$; number of shelves $p$; mold availabilities $\{c_j\}_{j\in J}$}
  \KwOut{A neighbor of $S$}
  \BlankLine
  \Fn{\TwoStep{$S$}}{
    Draw $(\ell^{*},r^{*})$ uniformly from $\Sigma(S)$; let $j^{*}$ be the job of $S_{\ell^{*}}[r^{*}]$
    \tcp*{switch only}
    \eIf{$c_{j^{*}} = 1$}{
      \Return \Relocate{\Relocate{$S$}} \tcp*{two independent relocations}
    }{
      \Return \Redistribute{$S$} \tcp*{single redistribution}
    }
  }
  \BlankLine
  \Fn{\Relocate{$S$}}{
    $S' \gets$ copy of $S$\;
    Draw $(\ell_s,r_s)$ uniformly from $\Sigma(S')$; $\sigma \gets S'_{\ell_s}[r_s]$
    \tcp*{quantity of $\sigma$ unchanged}
    Draw $\ell_t$ uniformly from $\{1,\dots,p\}$ \tcp*{$\ell_t=\ell_s$ allowed}
    Draw $r_t$ uniformly from $\{1,\dots,|S'_{\ell_t}|+1\}$ \tcp*{indexed before removal}
    Delete $\sigma$ from $S'_{\ell_s}$ at position $r_s$\;
    \lIf{$\ell_t=\ell_s$ \textbf{and} $r_t>r_s$}{$r_t \gets r_t-1$}
    Insert $\sigma$ into $S'_{\ell_t}$ at position $r_t$; \Return $S'$\;
  }
  \BlankLine
  \Fn{\Redistribute{$S$}}{
    $S' \gets$ copy of $S$;\quad $M \gets \{\, j \in J : c_j>1 \,\}$\;
    \lIf{$M=\emptyset$}{\Return $S'$}
    Draw $j^{\dagger}$ uniformly from $M$ \tcp*{independent of $j^{*}$}
    $P \gets$ positions of the $c_{j^{\dagger}}$ sub-jobs of $j^{\dagger}$, sorted by mold index\;
    $(q_1,\dots,q_{c_{j^{\dagger}}}) \gets{}$ \Split{$n_{j^{\dagger}}$, $c_{j^{\dagger}}$}
    \tcp*{uniform on $\{q\geq 0:\sum_i q_i=n_{j^{\dagger}}\}$}
    \lFor{$i=1$ \KwTo $c_{j^{\dagger}}$}{
      set the quantity of the $i$-th sub-job of $P$ to $q_i$ \tcp*{positions unchanged}
    }
    \Return $S'$\;
  }
  \BlankLine
  \Fn{\Split{$N,k$}}{
    \lIf{$k=1$ \textbf{or} $N=0$}{\Return $(N,0,\dots,0)$}
    Draw $\{d_1<\dots<d_{k-1}\}$ uniformly among the $(k-1)$-subsets of $\{1,\dots,N+k-1\}$\;
    $d_0 \gets 0$;\quad $d_k \gets N+k$;\quad \Return $(d_i-d_{i-1}-1)_{i=1}^{k}$\;
  }
\end{algorithm}

\section*{Correspondence with the implementation}

Algorithm~\ref{alg:neigh} is a line-by-line transcription of four functions of
the file \texttt{simulated\_annealing.jl}, in the directory \texttt{scripts} of
the project repository
\begin{center}
  \url{https://github.com/ndlopes-github/ShoeOptSetupTime}
\end{center}
\noindent
namely:

\begin{center}
\begin{tabular}{ll}
  \TwoStep{}     & \texttt{two\_step\_neighbor} \\
  \Relocate{}    & \texttt{neighbor\_ppartition} \\
  \Redistribute{} & \texttt{neighbor\_quantity\_redistribution} \\
  \Split{}       & \texttt{split\_quantity\_randomly} \\
\end{tabular}
\end{center}

\noindent
The initial solution is built by \texttt{create\_random\_partition} in the same
file.

\end{document}
